\section{Introduction}\label{sec:intro}

\subsection{The Problem}

Internet censorship research has traditionally focused on the network layer:
whether a website is DNS-poisoned, IP-blocked, or throttled somewhere between the
user and the server. A mature ecosystem of tools measures this activity, most
notably OONI, the Open Observatory of Network
Interference~\cite{filasto2012ooni}, and Censored
Planet~\cite{sundararaman2020censoredplanet}. But a growing amount of censorship
now happens \emph{above} the network, at the level of the mobile app stores
operated by Apple and Google. Because these two companies run a separate
storefront for each country, a government that wants to make an app disappear no
longer needs to touch the network at all. It can simply compel Apple or Google to
remove that app from its national store. For the ordinary smartphone user, an app
that is not in the store effectively does not exist, and no firewall or traffic
filtering is required to achieve that result.

The problem this project addresses is that we lack systematic, ongoing visibility
into this form of censorship. While the network layer is measured continuously by
OONI and others, the app-store layer is measured far less consistently, even
though it has become one of the most powerful chokepoints for controlling access
to VPNs, secure messaging, and independent news. We build a measurement pipeline
that queries the per-country storefronts of the Apple App Store and Google Play
for a fixed basket of censorship-sensitive apps, records which apps are available
in which countries, tracks changes over the term, and cross-validates the results
against transparency reports and documented events.

\subsection{Why It Matters}

App-store removals are neither hypothetical nor rare. They are a recurring,
government-driven mechanism with large real-world reach.

\paragraph{Russia (2024).} In July 2024, Apple removed at least 25 VPN apps from
the Russian App Store, including NordVPN, Proton VPN, and Private Internet Access,
after a demand from the state communications regulator
Roskomnadzor~\cite{techcrunch2024russia}. Apple's own notice to developers stated
that the apps would be removed ``per demand from Roskomnadzor'' because they
contained content ``illegal in Russia''~\cite{moscowtimes2024russia}. Independent
monitoring by the GreatFire App Censorship Project put the total at roughly 60
VPN apps removed over the summer of 2024~\cite{applecensorship}.

\paragraph{China.} Apple has removed VPN apps from the Chinese App Store since
2017 to comply with local law~\cite{cnbc2017china}. A review by the Tech
Transparency Project identified 3{,}257 apps missing from the China store, with
nearly a third relating to human-rights-sensitive topics such as privacy tools,
Tibetan Buddhism, Hong Kong protests, and LGBTQ+
issues~\cite{ttp2021china}.

\paragraph{India (2020).} India banned TikTok and dozens of other Chinese apps in
June 2020, a figure that later grew to more than 500 apps, affecting an estimated
200 million TikTok users in the country and removing the apps from both major
stores~\cite{techcrunch2020india}.

The importance of the problem is reinforced by the fact that dedicated
civil-society and academic projects have grown up specifically to track it, which
signals both demand and an unfilled gap. GreatFire's AppleCensorship project runs
an ``App Store Monitor'' that crawls 155 country stores to test
availability~\cite{applecensorship}, and Ververis et al.\ introduced a
methodology for querying app-store search engines and public APIs across
countries, finding that apps focused on Tibetan or LGBTQ+ topics face
disproportionately high censorship~\cite{ververis2019shedding}. Our project
extends this line of work (Section~\ref{sec:related}): existing coverage is
either event-driven news reporting or focused narrowly on VPNs in Russia and
China, whereas we run a controlled, longitudinal, multi-category measurement and
pair it with an explicit treatment of the attribution problem.

\subsection{Who Can Act on the Results}

Our results support several concrete actions, each mapped to actors equipped to
take them. \textbf{Journalists and civil-society organizations} (GreatFire,
Freedom House, Access Now, EFF) can use a longitudinal availability dataset to
detect and publicize removals close to when they happen, rather than relying on
developers to notice their own delisting. \textbf{The platform companies
themselves} (Apple, Google) are the only parties who can restore an app or
improve disclosure; independent measurement creates external pressure and a check
on transparency reports that, as the China case shows, can understate the scope
of removals. \textbf{Policymakers and regulators} in democracies, for example
bodies implementing the EU Digital Services Act, can use evidence of compelled
removals to inform platform-transparency and takedown-reporting rules.
\textbf{Circumvention-tool developers} can use the data to see which categories
and regions are most aggressively targeted and prioritize alternative
distribution channels (sideloading, third-party stores, web apps). Finally,
\textbf{users and the research community} gain an open dataset and reproducible
method that complements network-level tools such as OONI at a layer they do not
currently cover.
